\documentclass[10pt,a4paper]{amsart}
\usepackage[margin=19mm]{geometry}
\usepackage{amsmath,amssymb,mathtools}
\usepackage[scr=rsfs,cal=euler]{mathalfa}
\usepackage{xfrac}
\usepackage[unicode,hidelinks,bookmarks=false]{hyperref}
\newcommand{\ie}{i.e.\ }
\newcommand{\cf}{cf.\ }
\newcommand{\resp}{respectively\ }
\newcommand{\Subsection}{\subsection}
\providecommand{\nobreakdash}{\nobreak\hbox{-}}
\setlength{\emergencystretch}{2em}
\multlinegap0pt
\usepackage{amssymb}
\usepackage{xcolor}
\usepackage{float}
\usepackage{tikz}
\usepackage{mathtools}
\usepackage[shortlabels]{enumitem}
\newtheorem{lemma}{Lemma}
\newcommand{\F}{\mathfrak{F}}
\renewcommand{\L}{\mathcal{L}}
\title{Dual Cheeger constant for weighted graphs over ordered fields}
\author{Anna Muranova}
\date{Source: arXiv:2211.01654v1 (CC BY 4.0)}
\begin{document}
\maketitle

\noindent\emph{Source and licence.} Dual Cheeger constant for weighted graphs over ordered fields. Version arXiv:2211.01654v1 (CC BY 4.0), distributed by arXiv under the Creative Commons Attribution 4.0 International licence, http://creativecommons.org/licenses/by/4.0/, as recorded in the arXiv licence metadata for this exact version and shown on the abstract page https://arxiv.org/abs/2211.01654v1. Full licence notice: \texttt{licenses/arxiv-2211.01654-CC-BY-4.0.txt}.\par\smallskip

\noindent\textit{Editorial scope.} Counted unit: the article's lemma lemma (source0.tex lines 345--365). The article's complete original proof of it is delivered with it (source0.tex lines 367--428, one \texttt{proof} environment of 62 lines). No mathematics is written, repaired or re-derived: the statement and the proof are the article's own lines, in the article's own notation, and no retained line is substituted or re-flowed.  No further source-local context is needed: every cross-reference of every retained line resolves inside the two delivered spans.  Nothing was omitted from any retained span, so the package carries no omission record.  No retained line contains a citation, so the wrapper carries no bibliography and no bibliography entry.  No retained line contains a figure environment, an image-inclusion command or drawing code, so no image is delivered and the package declares no asset dependency.  Editorial formatting only, in addition: an AMS \texttt{amsart} preamble that restates the article's own declarations and macros used by the retained lines, namely the environments \texttt{lemma} on separate counters, together with the standard packages those lines need.  The retained environments print in this document as Lemma 1 in order of appearance, whereas the article prints the same items with the numbering of its own full document.  The complete native source of the article as distributed in the arXiv e-print for this exact version is bundled unchanged as \texttt{sources/132966/source0.tex}.
\begin{lemma}
Let $(V,b)$ be a graph without isolated vertices and $f\in \F_V$ be such that $V^+_f\ne\emptyset$ and $V\setminus V^+_f\ne\emptyset$. Let
% is such that $b(V^+_f)\preceq b(V^-_f)$. If
\begin{equation*}
h(f)=\min_{\emptyset\ne S\subset V^+_f}\dfrac{b(S, V\setminus S)}{b(S)}
\end{equation*}
and 
$$
f_+(x)=\begin{cases}
f(x), x\in V^+_f,\\
0, \mbox{ otherwise},
\end{cases}
$$
then
\begin{equation*}
1-\sqrt{1-h^2(f)}\preceq \dfrac{\langle\L f_+, f_+\rangle}{\langle f_+, f_+\rangle} \preceq 1+\sqrt{1-h^2(f)}.
\end{equation*}


%If we order all the vertices of the graph as $x_1,x_2\dots\x_N$ such that $f(x_i)\preceq f(x_{i+1})$
\end{lemma}

\begin{proof}
Let us denote
\begin{equation*}
W:=\dfrac{\langle\L f_+, f_+\rangle}{\langle f_+, f_+\rangle}.
\end{equation*}
By Green formula we can rewrite $W$ as 
\begin{equation*}
W=\dfrac{\frac{1}{2}\sum_{x,y\in V}(f_+(x)-f_+(y))^2b(x,y)}{\sum_{x\in V}f_+^2(x)b(x)}.
\end{equation*}
Then multiplying the numerator and the denominator by $\sum_{x,y\in V}(f_+(x)+f_+(y))^2b(x,y)$ and using Cauchy-Schwarz inequality we get
\begin{equation}\label{W1}
W\succeq \dfrac{\frac{1}{2}\left(\sum_{x,y\in V}b(x,y)\left|f_+^2(x)-f_+^2(y)\right|\right)^2}{\left(\sum_{x\in V}f_+^2(x)b(x)\right)\left(\sum_{x,y\in V}(f_+(x)+f_+(y))^2b(x,y)\right)}.
\end{equation}
From the other hand,
\begin{align*}
&W\cdot {\sum_{x\in V}f_+^2(x)b(x)}=\langle\L f_+,f_+\rangle\\
&=2\sum_{x\in V}f_+^2(x)b(x)-\left(\sum_{x\in V}f_+^2(x)b(x)+\sum_{x,y\in V}f_+(x)f_+(y)b(x,y)\right)\\
&=2\sum_{x\in V}f_+^2(x)b(x)-\left(\sum_{x,y\in V}f_+(x)(f_+(x)+f_+(y))b(x,y)\right)
\\
&=2\sum_{x\in V}f_+^2(x)b(x)-\left(\sum_{x,y\in V}f_+(y)(f_+(x)+f_+(y))b(x,y)\right)\\
&=2\sum_{x\in V}f_+^2(x)b(x)-\frac12\left(\sum_{x,y\in V}(f_+(x)+f_+(y))^2b(x,y)\right),
\end{align*}
where in the last two lines we have switched notations for $x,y$ in the second sum and sum up two lines and divided by two in the last line. Using this in denominator of \eqref{W1} we get
\begin{equation}\label{W2}
W\succeq\dfrac{\frac{1}{2}\left(\sum_{x,y\in V}b(x,y)\left|f_+^2(x)-f_+^2(y)\right|\right)^2}{\left(\sum_{x\in V}f_+^2(x)b(x)\right)^2\left(4-2 W\right)}.
\end{equation}
Now our goal is to estimate the numerator of \eqref{W2}
$$
A:=\frac{1}{2}\left(\sum_{x,y\in V}b(x,y)\left|f_+^2(x)-f_+^2(y)\right|\right)^2
$$ 
in terms of $h(f)$.
Firstly, let us order all the vertices of the graph as $x_1,x_2\dots,x_N$ in the way that $f(x_i)\preceq f(x_{i+1})$. Further, let us denote 
$$
S_k=\{x_{k+1},\dots,x_N\}\subset~V\;\;\;\mbox{ and }\;\;\;n_0=\max_i\{i\mid x_i\not \in V^+_f, 0\le i<N \},
$$
i.e. for any $k\ge n_0$ we have $S_k\subset~V^+_f$.

Now we can rewrite $A$ as
\begin{align*}
A&=2\left(\sum_{i=1}^{N-1}\sum_{j=i+1}^N b(x_i,x_j)\left(f_+^2(x_j)-f_+^2(x_i)\right)\right)^2\\
&=2\left(\sum_{i=1}^{N-1}\sum_{j=i+1}^N b(x_i,x_j)\sum_{k=i}^{j-1}\left(f_+^2(x_{k+1})-f_+^2(x_k)\right)\right)^2\\
&=2\left(\sum_{i=1}^{N-1}\sum_{j=i+1}^N \sum_{k=i}^{j-1}b(x_i,x_j)\left(f_+^2(x_{k+1})-f_+^2(x_k)\right)\right)^2\\
&=[\mbox{change limits of summation for $k$ and $j$}]\\
&=2\left(\sum_{i=1}^{N-1}\sum_{k=i}^{N-1} \sum_{j=k+1}^{N}b(x_i,x_j)\left(f_+^2(x_{k+1})-f_+^2(x_k)\right)\right)^2\\
&=[\mbox{change limits of summation for $i$ and $k$}]\\
&=2\left(\sum_{k=1}^{N-1}\underbrace{\sum_{i=1}^{k} \sum_{j=k+1}^{N}b(x_i,x_j)}_{b(S_k,V\setminus S_k)}\left(f_+^2(x_{k+1})-f_+^2(x_k)\right)\right)^2\\
&=2\left(\sum_{k=1}^{N-1}{b(S_k,V\setminus S_k)}\left(f_+^2(x_{k+1})-f_+^2(x_k)\right)\right)^2\\
&=2\left(\sum_{k=n_0}^{N-1}{b(S_k,V\setminus S_k)}\left(f_+^2(x_{k+1})-f_+^2(x_k)\right)\right)^2\\
&\succeq 2\left(\sum_{k=n_0}^{N-1}{h(f)b(S_k)}\left(f_+^2(x_{k+1})-f_+^2(x_k)\right)\right)^2\\
&= 2h^2(f)\left(\sum_{k=n_0}^{N-1}{\sum_{i=k+1}^{N}b(x_i)}\left(f_+^2(x_{k+1})-f_+^2(x_k)\right)\right)^2\\
&=[\mbox{change limits of summation for $i$ and $k$}]\\
&=2h^2(f)\left(\sum_{i=n_0+1}^{N}{\sum_{k=n_0}^{i-1}b(x_i)}\left(f_+^2(x_{k+1})-f_+^2(x_k)\right)\right)^2\\
&=2h^2(f)\left(\sum_{i=n_0+1}^{N}{b(x_i)}\left(f_+^2(x_i)-f_+^2(x_{n_0})\right)\right)^2\\
&=2h^2(f)\left(\sum_{i=n_0+1}^{N}{b(x_i)}f_+^2(x_i)\right)^2=2h^2(f)\left(\sum_{x\in V}{b(x)}f_+^2(x)\right)^2,
\end{align*}
where we have used several times that $f_+(x_k)=0$ for any $k\le n_0$.
Applying the last estimate to \eqref{W2} we obtain 
$$
W \succeq \dfrac{2 h^2(f)}{(4-2W)},
$$
from where immediately follows the statement of the lemma.
\end{proof}

\end{document}
